\subsection{Concurrent Knowledge Base}
\label{sec:ckb}

Based on the extracted concurrency entities and interaction DSL described in the previous section, we construct a \emph{Concurrent Knowledge Base (CKB)} to organize concurrency-relevant information in a hierarchical and queryable form.
The CKB serves as the structured knowledge layer that supports hierarchical retrieval and LLM-guided fault localization.

Formally, the Concurrent Knowledge Base is defined as:
\[
\mathcal{K} = \{ \mathcal{C}, \mathcal{P}, \mathcal{E}, \mathcal{M}, \mathcal{I} \}
\]
where each element captures concurrency knowledge at a different abstraction level.

The Concurrent Knowledge Base (CKB) organizes concurrency-relevant information at multiple abstraction levels to support progressive fault localization.
At coarse granularity, component- and package-level knowledge captures high-level system structure and functionality.
Each component and package is associated with brief descriptions summarizing its responsibilities and concurrency relevance, derived from build configurations, directory structure, and developer documentation.
This information provides early-stage guidance for narrowing the search space when analyzing large codebases.

At finer granularity, the CKB records detailed concurrency semantics at the thread-entry, method, and interaction levels.
Formally, thread-entry knowledge corresponds to the set $\mathbb{T}_e$, where each $t \in \mathbb{T}_e$ represents a concurrency entry point such as a thread, executor submission, or asynchronous callback.
For each thread entry $t$, the CKB records the set of other entries $t' \in \mathbb{T}_e$ with which $t$ may interact through shared resource accesses.
Method-level knowledge is defined over the set of methods $\mathbb{M}$.
For each method $m \in \mathbb{M}$, the CKB records whether $m$ serves as a thread entry ($m \in \mathbb{T}_e$), the set of shared variables $v \in \mathbb{V}_s$ accessed by $m$, and the corresponding access types (e.g., read or write).
This information captures how concurrency behavior manifests at the code level.
Finally, interaction-level knowledge consists of interaction DSL elements $\mathcal{I}$.
Each interaction $I \in \mathcal{I}$ summarizes a potential concurrent behavior between a pair of distinct thread entries $(t_1, t_2)$, where $t_1, t_2 \in \mathbb{T}_e$.

\vspace{-0.5em}
\begin{table}[h]
\centering
\caption{Function calls for hierarchical retrieval}
\vspace{-8pt}
\label{tab:function_calls}
\scriptsize
\begin{tabular}{l|l|p{7.5cm}}
\hline
\textbf{Name} & \textbf{Arguments} & \textbf{Description} \\
\hline
list\_components & None &
Return all components with brief descriptions. \\
\hline
list\_packages\_under\_component & component\_name &
Return all packages under the specified component. \\
\hline
list\_entries\_under\_package & package\_name &
Return all thread entry points defined in the package. \\
\hline
list\_method\_under\_class & class\_name &
Return all methods under the specified class. \\
\hline
get\_entry\_summary & entry\_name &
Return concurrency summary of the thread entry and its interacting entries. \\
\hline
get\_method\_summary & method\_name &
Return concurrency-related information of the method. \\
\hline
get\_interaction\_dsl & entry\_pair &
Return interaction DSL for the given pair of thread entries. \\
\hline
get\_code\_snippet & method\_name &
Return code snippet of the specified method. \\
\hline
exit & None &
Terminate the retrieval process. \\
\hline
\end{tabular}
\end{table}
\vspace{-8pt}
\vspace{-0.5em}

\paragraph{Retrieval Functions}
\label{subsec:ckb-retrieval}

To enable hierarchical and controlled access to the Concurrent Knowledge Base, we expose a small set of retrieval functions, summarized in Table~\ref{tab:function_calls}.
Each function corresponds to a specific abstraction level in the knowledge base, ranging from components and packages to thread entries, methods, and interaction DSL elements.
By restricting retrieval to these predefined functions, the framework ensures that the LLM incrementally explores only concurrency-relevant context, preventing uncontrolled traversal of the codebase and keeping the reasoning process focused and interpretable.
